% GENERATED FILE. Edit canonical lessons, then run npm run generate:books.
% canonical-source-hash: fnv1a64:17b9de0e435209b6
% canonical-lessons: ZH-C460-yumen, ZH-C460-tanqi, ZH-C460-nanweiqing, ZH-C460-tongqing, ZH-C460-xianmu

\chapter{Feeling Low}
\label{ch:zh-feeling-low}
\hlchaptermodality{460}

\begin{chapteropening}
\textbf{By the end of this chapter:} \emph{I can say I feel down, ask why someone sighs, say I am embarrassed and that I feel for someone or envy them.}

Complete the last lesson of chapter 460: I can say I feel down, ask why someone sighs, say I am embarrassed and that I feel for someone or envy them.
\end{chapteropening}

\section[yùmèn]{yùmèn — down; gloomy}

\label{lesson:ZH-C460-yumen}



\begin{quote}
Before the new one: say the Mandarin for what is really on one's mind, then the Mandarin for hurt by unfair blame.
\end{quote}

\subsection*{You'll want to know: yùmèn}
\emph{yùmèn} — \textquotedblleft{}down; gloomy\textquotedblright{}.

\begin{quote}
\emph{wǒ bìng le yí ge xīngqī hěn yùmèn} — I've been ill for a week, and I'm feeling really down.
\end{quote}

Its characters come later; for now it is a word for the ear.

\begin{cousinweb}
\emph{yùmèn} is built from \emph{yù} 'dense' and \emph{mèn} 'vexed', both inherited from Middle Chinese, \emph{'jut} and \emph{mwonH} (Guangyun, in Baxter's transcription), which glosses the first as thick and blocked and the second 'vexed'.
\end{cousinweb}

\begin{grammarlens}[title={What you've built}]
One more word for saying how it was.
\end{grammarlens}

\subsection*{Your turn}
\begin{itemize}
\raggedright
  \item \emph{Say it:} \emph{yùmèn}
  \item \emph{Say it:} \emph{yùmèn}, once more
  \item \emph{Say it:} \emph{wǒ bìng le yí ge xīngqī hěn yùmèn} — I've been ill for a week, and I'm feeling really down.
  \item \emph{Recall:} say the Mandarin for what is really on one's mind, then the Mandarin for hurt by unfair blame, then say \emph{yùmèn} again
\end{itemize}

\begin{tcolorbox}[breakable,colback=teal!4,colframe=teal!35!black,title={Before you move on}]
What does yùmèn mean? (\textquotedblleft{}Down; gloomy\textquotedblright{}.) Say it once more.
\end{tcolorbox}
\section[tànqì]{tànqì — to sigh}

\label{lesson:ZH-C460-tanqi}



\begin{quote}
Before the new one: say the Mandarin for hurt by unfair blame, then the Mandarin for down in the dumps.
\end{quote}

\subsection*{You'll want to know: tànqì}
\emph{tànqì} — \textquotedblleft{}to sigh\textquotedblright{}.

\begin{quote}
\emph{nǐ wèishénme tànqì} — Why are you sighing?
\end{quote}

Its characters come later; for now it is a word for the ear.

\begin{cousinweb}
\emph{tànqì} is built from \emph{tàn} 'to sigh' and \zh{气} 'breath', both inherited from Middle Chinese, \emph{thanH} and \emph{khj+jH} (Guangyun, in Baxter's transcription), which glosses the first as a long breath and the second 'breath'.
\end{cousinweb}

\begin{grammarlens}[title={What you've built}]
One more word for saying how it was.
\end{grammarlens}

\subsection*{Your turn}
\begin{itemize}
\raggedright
  \item \emph{Say it:} \emph{tànqì}
  \item \emph{Say it:} \emph{tànqì}, once more
  \item \emph{Say it:} \emph{nǐ wèishénme tànqì} — Why are you sighing?
  \item \emph{Recall:} say the Mandarin for hurt by unfair blame, then the Mandarin for down in the dumps, then say \emph{tànqì} again
\end{itemize}

\begin{tcolorbox}[breakable,colback=teal!4,colframe=teal!35!black,title={Before you move on}]
What does tànqì mean? (\textquotedblleft{}To sigh\textquotedblright{}.) Say it once more.
\end{tcolorbox}
\section[\texorpdfstring{\zh{难为情} (nánwéiqíng)}{难为情 (nánwéiqíng)}]{\zh{难为情} (nánwéiqíng) — embarrassed}

\label{lesson:ZH-C460-nanweiqing}



\begin{quote}
Before the new one: say the Mandarin for down in the dumps, then the Mandarin for to sigh.
\end{quote}

\subsection*{You'll want to know: \zh{难为情}}
\textbf{\zh{难为情}} — \emph{nánwéiqíng} — \textquotedblleft{}embarrassed\textquotedblright{}.

\begin{quote}
\textbf{\zh{大家都笑我我很难为情}} \emph{dàjiā dōu xiào wǒ wǒ hěn nánwéiqíng} — Everyone laughed at me, and I was so embarrassed.
\end{quote}

\begin{cousinweb}
\textbf{\zh{难为情}} (\emph{nánwéiqíng}) is built from \zh{难} 'hard', \zh{为} 'to make' and \zh{情} 'feeling', all inherited from Middle Chinese, \emph{nan}, \emph{hjwe} and \emph{dzjeng} (Guangyun, in Baxter's transcription). The dictionary glosses the first 'hard to do'; the sense 'embarrassed' is not given there.
\end{cousinweb}

\begin{grammarlens}[title={What you've built}]
One more word for saying how it was.
\end{grammarlens}

\subsection*{Your turn}
\begin{itemize}
\raggedright
  \item \emph{Say it:} \emph{nánwéiqíng}
  \item \emph{Say it:} \emph{nánwéiqíng}, once more
  \item \emph{Say it:} \emph{dàjiā dōu xiào wǒ wǒ hěn nánwéiqíng} — Everyone laughed at me, and I was so embarrassed.
  \item \emph{Recall:} say the Mandarin for down in the dumps, then the Mandarin for to sigh, then say \emph{nánwéiqíng} again
\end{itemize}

\begin{tcolorbox}[breakable,colback=teal!4,colframe=teal!35!black,title={Before you move on}]
What does \zh{难为情} mean? (\textquotedblleft{}Embarrassed\textquotedblright{}.) Say it once more.
\end{tcolorbox}
\section[\texorpdfstring{\zh{同情} (tóngqíng)}{同情 (tóngqíng)}]{\zh{同情} (tóngqíng) — to sympathize with}

\label{lesson:ZH-C460-tongqing}



\begin{quote}
Before the new one: say the Mandarin for to sigh, then the Mandarin for embarrassed.
\end{quote}

\subsection*{You'll want to know: \zh{同情}}
\textbf{\zh{同情}} — \emph{tóngqíng} — \textquotedblleft{}to sympathize with\textquotedblright{}.

\begin{quote}
\emph{nàge háizi gēn māma zǒusàn le wǒ hěn tóngqíng tā} — That boy got separated from his mother, and I really feel for him.
\end{quote}

\begin{cousinweb}
\textbf{\zh{同情}} (\emph{tóngqíng}) is built from \zh{同} 'the same' and \zh{情} 'feeling', both inherited from Middle Chinese, \emph{duwng} and \emph{dzjeng} (Guangyun, in Baxter's transcription), which glosses the first 'equal' and 'together'.
\end{cousinweb}

\begin{grammarlens}[title={What you've built}]
One more word for saying how it was.
\end{grammarlens}

\subsection*{Your turn}
\begin{itemize}
\raggedright
  \item \emph{Say it:} \emph{tóngqíng}
  \item \emph{Say it:} \emph{tóngqíng}, once more
  \item \emph{Say it:} \emph{nàge háizi gēn māma zǒusàn le wǒ hěn tóngqíng tā} — That boy got separated from his mother, and I really feel for him.
  \item \emph{Recall:} say the Mandarin for to sigh, then the Mandarin for embarrassed, then say \emph{tóngqíng} again
\end{itemize}

\begin{tcolorbox}[breakable,colback=teal!4,colframe=teal!35!black,title={Before you move on}]
What does \zh{同情} mean? (\textquotedblleft{}To sympathize with\textquotedblright{}.) Say it once more.
\end{tcolorbox}
\section[xiànmù]{xiànmù — to envy}

\label{lesson:ZH-C460-xianmu}



\begin{quote}
Before the new one: say the Mandarin for embarrassed, then the Mandarin for to sympathize with.
\end{quote}

\subsection*{You'll want to know: xiànmù}
\emph{xiànmù} — \textquotedblleft{}to envy\textquotedblright{}.

\begin{quote}
\emph{nǐ qù guo hěn duō guójiā wǒ hěn xiànmù nǐ} — You've been to so many countries; I really envy you.
\end{quote}

Its characters come later; for now it is a word for the ear.

\begin{cousinweb}
\emph{xiànmù} is built from \emph{xiàn} 'to covet' and \emph{mù} 'to long for', both inherited from Middle Chinese, \emph{zjenH} and \emph{muH} (Guangyun, in Baxter's transcription). The dictionary lists the first under three readings and glosses this one 'to covet'; it glosses the second as thinking of someone with longing.
\end{cousinweb}

\begin{grammarlens}[title={What you've built}]
One more word for saying how it was.
\end{grammarlens}

\subsection*{Your turn}
\begin{itemize}
\raggedright
  \item \emph{Say it:} \emph{xiànmù}
  \item \emph{Say it:} \emph{xiànmù}, once more
  \item \emph{Say it:} \emph{nǐ qù guo hěn duō guójiā wǒ hěn xiànmù nǐ} — You've been to so many countries; I really envy you.
  \item \emph{Recall:} say the Mandarin for embarrassed, then the Mandarin for to sympathize with, then say \emph{xiànmù} again
\end{itemize}

\begin{tcolorbox}[breakable,colback=teal!4,colframe=teal!35!black,title={Before you move on}]
What does xiànmù mean? (\textquotedblleft{}To envy\textquotedblright{}.) Say it once more.
\end{tcolorbox}
